%% Marketing Letters -- Replication Corner
%% Manuscript: Does AI Reduce the Reputation Premium in Knowledge Diffusion?
%%             A Conceptual Replication Using Chinese Journal Articles

\documentclass[12pt,a4paper]{article}

% ---- Packages ---------------------------------------------------------------
\usepackage[top=2.54cm,bottom=2.54cm,left=2.54cm,right=2.54cm]{geometry}
\usepackage{setspace}
\usepackage{amsmath,amssymb}
\usepackage{booktabs}
\usepackage{threeparttable}
\usepackage{multirow}
\usepackage{array}
\usepackage[authoryear,sort]{natbib}
\usepackage[hidelinks]{hyperref}
\usepackage{microtype}
\usepackage{times}
\usepackage[utf8]{inputenc}
\usepackage[T1]{fontenc}
\usepackage{caption}
\captionsetup{labelfont=bf, font=small}
\newcommand{\doi}[1]{\href{https://doi.org/#1}{doi:#1}}

\doublespacing

% ---- Title block ------------------------------------------------------------
\title{\textbf{Does AI Reduce the Reputation Premium\\
in Knowledge Diffusion?\\
A Conceptual Replication Using Chinese Journal Articles}\\[0.8em]
\large\textit{Replication Corner / Marketing Letters}}

\author{[Blinded for review]}

\date{}

\begin{document}

\maketitle

% ---- Abstract ---------------------------------------------------------------
\begin{abstract}
\noindent
\citet{JoLi2026} report that AI-style writing attenuates the
institutional reputation premium in SSRN marketing working papers.
We conceptually replicate this finding using 843 peer-reviewed Chinese
business and management journal articles indexed on CNKI
(China National Knowledge Infrastructure; 2023--2026).
AI-like structural writing style is proxied by an abstract-level
dictionary score quantifying the density of 23 pre-specified
AI-associated writing markers; institutional prestige is coded on a
four-tier classification based on China's C9, 985, and
211/Double First-Class university designations.
The AI-style proxy independently predicts greater download volume
($\hat{\beta}=0.22$, $p<.001$), partially replicating \citet{JoLi2026}.
However, institutional prestige is not a significant predictor of
downloads or citations in any specification (all $p>.10$), consistent
with editorial certification reducing the marginal role of institutional
prestige in formal publication venues---a pattern not observed on SSRN.
The Prestige\,$\times$\,AI-style interaction is near zero and
statistically insignificant across all four specifications (binary and
continuous prestige; download and citation outcomes).
We do not find evidence that AI-style writing attenuates a reputation
premium in this peer-reviewed context; more fundamentally, the prestige
premium itself appears absent, suggesting a boundary condition for the
mechanism documented by \citet{JoLi2026}: AI can attenuate
reputation-based diffusion only where a reputation premium is operative.

\smallskip\noindent
\textbf{Keywords:} AI writing style; institutional prestige; knowledge
diffusion; reputation premium; replication; CNKI; China
\end{abstract}

\newpage

% ---- 1. Introduction --------------------------------------------------------
\section{Introduction}
\label{sec:intro}

Generative AI tools have transformed academic writing, raising a
consequential question for the sociology of knowledge: does AI-assisted
writing democratize scholarly attention by reducing the visibility
advantage of researchers from elite institutions?
If AI helps less-established authors produce work that is
indistinguishable in presentation quality from elite-institution output,
institutional affiliation may become a weaker quality heuristic for
readers, and the longstanding reputation premium in academic knowledge
diffusion may narrow
\citep{Merton1968,Spence1973,DiMaggio1983}.

\citet{JoLi2026} provide the first direct evidence on this question,
using a sample of SSRN marketing working papers.
They find that articles with higher AI-style writing scores attract
more downloads and readership, and---critically---that AI-style writing
moderates the institutional prestige premium: the positive association
between elite-university affiliation and diffusion is attenuated when
the paper exhibits AI-like writing characteristics.
The authors interpret this as evidence that AI is beginning to level
the playing field in academic knowledge dissemination, at least in the
preprint ecosystem.

We examine whether this finding generalizes to a structurally distinct
context: peer-reviewed Chinese business and management journal articles
indexed on CNKI.
Two key features distinguish this setting from SSRN.
First, CNKI journals impose editorial gatekeeping---peer review and
editorial selection---that SSRN preprints lack entirely.
This gatekeeping provides an independent quality signal correlated with
institutional prestige, potentially preserving the prestige premium
even when writing style converges \citep{Bornmann2008}.
Second, the Chinese academic system embeds institutional hierarchy
deeply in government funding allocation, faculty promotion criteria,
and journal-composition norms, suggesting that prestige operates
through channels beyond presentation quality that AI-style writing
cannot substitute \citep{DiMaggio1983,Long1979}.
For both reasons, we expect the moderating effect of AI-style writing
to be weaker or absent in this setting.

We test these expectations using 843 CNKI articles published between
2023 and 2026 on topics spanning AI, enterprise management, marketing,
innovation, ESG, and supply chain.
AI-like writing style is proxied by an abstract-level dictionary score
that counts the presence of AI-associated structural markers in each
Chinese abstract.
Institutional prestige is captured via a four-tier classification
based on China's C9, 985, and 211/Double First-Class university
designations.
Downloads and citations serve as diffusion outcomes.
We find that AI-style writing predicts greater download-based
diffusion---partially replicating \citet{JoLi2026}---but institutional
prestige is not significantly associated with diffusion in this sample.
The null prestige main effect is itself a boundary condition with
implications discussed in Section~\ref{sec:discussion}.
The critical interaction is consistently small and statistically
insignificant across all specifications.

Section~\ref{sec:theory} reviews the theoretical background and states
our propositions.
Section~\ref{sec:method} describes data and method.
Section~\ref{sec:results} reports results.
Section~\ref{sec:discussion} discusses boundary conditions and limitations.

% ---- 2. Theory and Hypotheses -----------------------------------------------
\section{Theoretical Background and Hypotheses}
\label{sec:theory}

\subsection{The Institutional Reputation Premium}

Signaling theory holds that observable attributes function as proxies
for unobservable quality under information asymmetry \citep{Spence1973}.
In academic knowledge diffusion, institutional affiliation serves as a
credible quality signal because elite institutions select, train, and
resource researchers in ways correlated with output quality
\citep{Long1979,Lariviere2010}.
Institutional theory further suggests that organizational hierarchy
becomes self-reinforcing through legitimating processes: readers,
reviewers, and editors allocate attention partly on the basis of
normalized prestige categories \citep{DiMaggio1983}.
Consistent with this logic, empirical research documents robust
citation and readership premiums for articles from high-prestige
institutions, controlling for article content \citep{Wang2013}.
\begin{description}
  \item[H1] \textit{(Prestige--Diffusion Association):}
  Articles affiliated with higher-prestige institutions will receive
  more downloads and more citations on CNKI than articles from
  lower-prestige institutions, controlling for publication age
  and article characteristics.
\end{description}
We note, however, that the peer-reviewed CNKI context may attenuate
this premium relative to the SSRN preprint setting; we treat H1 as an
empirical replication test rather than a strong theoretical prediction,
and interpret null results as substantively informative rather than
merely inconclusive.

\subsection{AI-Style Writing and Diffusion}

The proliferation of large language models since late 2022 has produced
a recognizable stylistic pattern in academic abstracts: structured
section-by-section organization, high density of methodological keywords,
formulaic transitions, and templated summary phrases
\citep{Liang2024,Biber2016}.
\citet{Liang2024} document the increasing prevalence and characteristic
patterns of LLM-assisted academic writing.
Although \citet{Biber2016} focus on English-language academic prose,
the structural properties that distinguish their register---formulaic
section organization, templated summary phrases, stereotyped transitional
markers---are observable in parallel form in Chinese LLM-assisted
abstracts, a cross-linguistic parallel that grounds our marker taxonomy.
Readers may interpret this style as signaling methodological
thoroughness and genre conformity, independent of substantive quality
differences.
\citet{JoLi2026} show that higher AI-style scores on SSRN predict
greater readership.
We test whether this association replicates in a peer-reviewed context:
\begin{description}
  \item[H2] \textit{(AI-Style--Diffusion Association):}
  Articles exhibiting higher AI-style writing scores will receive
  more downloads and more citations on CNKI, controlling for
  institutional prestige, publication age, and abstract
  characteristics.
\end{description}

\subsection{The Moderation Hypothesis and a Boundary Condition}

\citet{JoLi2026} argue that AI-style writing attenuates the prestige
premium by homogenizing presentation quality: when AI makes elite and
non-elite authors stylistically indistinguishable, readers rely less
on institutional affiliation as a quality heuristic.
This mechanism presupposes that writing style constitutes the primary
channel through which prestige influences diffusion.
In the SSRN preprint ecosystem---where no editorial gatekeeping
occurs---institutional affiliation and writing quality are indeed
among the few observable signals available to readers, making this
substitution plausible.

In the CNKI journal context, a third signal is available: peer review
certification.
Journal acceptance signals that independent reviewers have validated
the article's content \citep{Bornmann2008}.
This certification is correlated with institutional prestige---elite
institutions may have greater access to skilled reviewers and stronger
editorial networks---but is not reducible to writing style.
AI-style writing cannot crowd out the prestige premium when prestige
also operates through a certified-quality channel.
Furthermore, in the Chinese academic system, prestige classifications
(C9, 985, 211) are formally institutionalized in government funding
formulas and promotion criteria, amplifying the salience of
institutional affiliation for readers and editors alike.
Notably, high-prestige institutions in our sample produce abstracts
with marginally \textit{lower} AI-style scores (mean\,=\,0.048)
than low-prestige institutions (mean\,=\,0.050), which further
limits the scope for AI-style writing to equalize stylistic quality
across institutional tiers.
Based on this boundary-condition reasoning, we tentatively predict
that the Prestige\,$\times$\,AI-style interaction will be weaker in
magnitude and less consistently negative in the CNKI peer-reviewed
context than in the SSRN preprint context of \citet{JoLi2026}.
We nevertheless frame this as an open research question because the
direction and magnitude of any attenuation depend on empirical
parameters that cannot be fixed ex~ante:
\begin{description}
  \item[RQ1] \textit{(Moderation by AI-Style Writing):}
  Does AI-style writing score moderate (attenuate) the positive
  association between institutional prestige and article diffusion
  outcomes on CNKI, and---if so---is the interaction weaker than
  the effect documented for SSRN preprints?
\end{description}

% ---- 3. Data and Method -----------------------------------------------------
\section{Data and Method}
\label{sec:method}

\subsection{Sample}

We searched CNKI for Chinese-language journal articles in the
economics, business, and management categories published between
January 2023 and June 2026, using keyword sets covering
AI/generative AI, digital transformation, corporate governance,
supply chain, ESG, marketing, and innovation.
We initially collected 847 CNKI records; four records lacked formal
publication identifiers or finalized publication information and were
excluded, yielding a final analytic sample of \textbf{843 articles}.
All articles come from CNKI journals in business-relevant categories
(economics, management, and business administration), corresponding
to the business and management scope of \citet{JoLi2026}'s SSRN sample.
Topic domains include AI and digital transformation, supply chain and
operations, corporate governance, ESG, and marketing and innovation.
Download and citation counts were retrieved from CNKI in June 2026,
providing a cross-sectional snapshot of article-level diffusion.

\subsection{Outcome Variables}

We use two outcomes: (1) $\log(\text{downloads}+1)$ and
(2) $\log(\text{citations}+1)$.
The log-plus-one transformation accommodates zero counts prevalent
among recently published articles, where many have yet to accumulate
citations.
To address left-truncation in the citation outcome most directly,
the robustness specification that excludes 2026 articles
(Table~\ref{tab:robustness}, Specification~3) restricts the sample
to articles with at least six months of potential citation exposure.

\subsection{AI-Style Score}

The \textit{ai\_style\_score} is a marker-based AI-style index,
constructed by scoring the binary presence of each of
23~pre-specified AI writing-style markers in the Chinese abstract text
and dividing the matched-marker count by 23, yielding a $[0,1]$
score.\footnote{The complete marker list with Chinese original text,
  English translation, and example abstracts is provided in
  Web Appendix~A2.}
The score ranges from 0 to 0.261 (mean\,=\,0.049, SD\,=\,0.050,
median\,=\,0.044), indicating that most abstracts contain at most
one or two matched markers.
The marker dictionary was constructed in two stages.
First, we adapted the AI writing-style taxonomy of \citet{Liang2024}
to Chinese academic social science by reviewing a development sample
of 2023--2024 abstracts classified by two bilingual coders; final
markers include structured section-phrases
(e.g., \textit{jizhi fenxi} [`mechanism analysis'],
\textit{yizhixing fenxi} [`heterogeneity analysis'])
and formulaic template phrases.
Inter-coder agreement on the development sample was
$\kappa = 0.79$, exceeding the 0.70 acceptability threshold.
Second, for each target abstract, we assessed the binary presence of
each marker via exact string matching; the score equals the proportion
of markers present.
This fully deterministic scoring procedure requires no external API
and is entirely reproducible from the published marker list.
For regressions the score is standardized (\textit{ai\_z}); a binary
median-split indicator (\textit{ai\_high}) is used in robustness checks.
The marker dictionary was calibrated on 2023--2024 writing conventions;
potential stylistic drift in 2025--2026 abstracts represents a
construct validity limitation addressed further in Section~\ref{sec:discussion}.

\subsection{Institutional Prestige}

Institutional prestige is coded from author affiliation fields on a
four-tier ordinal scale: 1\,=\,no nationally ranked institution;
2\,=\,211 Project or Double First-Class (DFC) institution;
3\,=\,985 Project institution; 4\,=\,C9 elite university.
Each article's score equals the maximum tier across all affiliated
institutions.
We adopt the maximum-tier rule---rather than first-author or
corresponding-author affiliation---because Chinese journal articles
list all co-authors' institutions in the byline, making the
highest-prestige affiliation the most visible signal to readers;
sensitivity analyses using first-author institution are reported in
Web Appendix~A3, and yield substantively identical conclusions.
Prestige coding was performed by two trained research assistants;
inter-rater agreement on a 40-article subsample was
$\kappa = 0.93$, indicating near-perfect reliability.
For primary analyses we use a binary indicator
(\textit{Prestige\_High}\,=\,1 if the prestige score\,$\geq 2$),
yielding 563 high-prestige articles (66.8\%) and 280 low-prestige
articles (33.2\%).
In robustness checks we standardize the four-tier score
(\textit{prestige\_score\_z}).

\subsection{Controls}

We control for $\log(\text{Age})$---the log of publication age in
days from publication to data collection (June 2026), capturing
cumulative exposure time---and $\log(\text{abstract length})$,
which controls for information density in the abstract.
Year fixed effects (2023\,=\,reference) absorb platform-level growth
trends and secular increases in AI-style writing over the sample period.

\subsection{Estimation}

We estimate OLS with HC3 heteroskedasticity-robust standard errors
\citep{MacKinnon1985}.
Two primary specifications are estimated for each outcome:

\noindent\textbf{Equation (1) --- Binary prestige:}
\begin{align}
\log(Y_i+1) &= \alpha
  + \beta_1\,\textit{Prestige\_High}_i
  + \beta_2\,\textit{ai\_z}_i
  + \beta_3\bigl(\textit{Prestige\_High}_i\times\textit{ai\_z}_i\bigr)
  \notag\\
  &\quad
  + \beta_4\log(\textit{Age}_i)
  + \beta_5\log(\textit{AbsLength}_i)
  + \textstyle\sum_{t}\delta_t\,\textit{Year}_t
  + \varepsilon_i
\label{eq:binary}
\end{align}

\noindent\textbf{Equation (2) --- Continuous prestige score:}
\begin{align}
\log(Y_i+1) &= \alpha
  + \beta_1\,\textit{prestige\_score\_z}_i
  + \beta_2\,\textit{ai\_z}_i
  + \beta_3\bigl(\textit{prestige\_score\_z}_i\times\textit{ai\_z}_i\bigr)
  \notag\\
  &\quad
  + \beta_4\log(\textit{Age}_i)
  + \beta_5\log(\textit{AbsLength}_i)
  + \textstyle\sum_{t}\delta_t\,\textit{Year}_t
  + \varepsilon_i
\label{eq:score}
\end{align}

\noindent
where $Y_i\in\{\text{downloads}_i,\,\text{citations}_i\}$.
The coefficient $\beta_3$ is the parameter of central inferential
interest; we report its point estimate, HC3 standard error, and 95\%
confidence interval.
Negative-binomial regressions on raw counts are reported in the
Web Appendix as a robustness check.

% ---- 4. Results -------------------------------------------------------------
\section{Results}
\label{sec:results}

\subsection{Descriptive Statistics}

Table~\ref{tab:descriptives} presents descriptive statistics and
bivariate correlations.
The mean AI-style score is 0.049 (SD\,=\,0.050).
High-prestige articles exhibit a marginally lower mean AI-style score
(0.048) than low-prestige articles (0.050); the bivariate correlation
between \textit{Prestige\_High} and \textit{ai\_z} is $r\,=\,{-}.03$
($p\,=\,.407$), confirming the two focal predictors are not collinear.
At the time of data collection, approximately 25\% of articles ($n = 211$)
had received zero citations; 2026 articles have had minimal citation
exposure, making log(Age) an especially important control for citations.
The correlation between $\log(\text{Age})$ and log downloads is
$r\,=\,.56$; the substantially higher correlation with log citations
($r\,=\,.79$) reflects that citation accumulation is more steeply
age-dependent than download accumulation in this sample.

\subsection{H1: Prestige Main Effect}

Contrary to H1, institutional prestige is not statistically significantly
associated with either outcome across any specification
(Table~\ref{tab:main}).
In the binary-prestige download model (Column~1),
$\hat{\beta}_{\text{Prestige}} = 0.113$ (SE\,=\,0.069, $p\,=\,.104$);
the corresponding bivariate correlation is $r\,=\,.05$ ($p\,=\,.185$),
indicating that prestige carries negligible unconditional association
with download volume.
The prestige--citation association is similarly non-significant
($\hat{\beta} = 0.046$, SE\,=\,0.065, $p\,=\,.475$; Column~2).
Using the standardized four-tier prestige score (Columns 3--4),
the pattern is identical: $\hat{\beta} = 0.031$ ($p\,=\,.366$) for
downloads and $\hat{\beta} = 0.016$ ($p\,=\,.606$) for citations.
The null prestige main effect is discussed as a boundary condition in
Section~\ref{sec:discussion}.

\subsection{H2: AI-Style Writing Main Effect}

Supporting H2 for the download outcome, the AI-style score independently
predicts download-based diffusion across all specifications.
In the binary-prestige download model,
$\hat{\beta}_{\text{AI}} = 0.216$ (SE\,=\,0.048, $p<.001$),
indicating that a one-standard-deviation increase in AI-style score
is associated with approximately 0.22 additional log-download units,
corresponding to a 24\% increase in expected downloads.
In the continuous prestige-score model (Column~3),
$\hat{\beta}_{\text{AI}} = 0.187$ (SE\,=\,0.032, $p<.001$;
point estimate slightly lower than Column~1, with a tighter SE
reflecting the switch to continuous-prestige specification).
The AI-style--citation association is positive but does not reach
conventional significance thresholds in the binary-prestige model
($\hat{\beta} = 0.059$, SE\,=\,0.043, $p\,=\,.178$; Column~2),
and is marginally significant in the continuous prestige model
($\hat{\beta} = 0.065$, SE\,=\,0.028, $p\,=\,.022$; Column~4).
The download-based replication of H2 is robust to the choice of
prestige operationalization.

\subsection{RQ1: The Prestige\,$\times$\,AI-Style Interaction}

The key test concerns $\hat{\beta}_3$, the interaction between prestige
and AI-style writing.
Across all four specifications, $\hat{\beta}_3$ is consistently small
and statistically insignificant (Table~\ref{tab:main}).
In the binary-prestige download model (Column~1),
$\hat{\beta}_3 = -0.044$ (SE\,=\,0.062, $p\,=\,.477$;
95\,\%\,CI: [$-$0.164, $+$0.077]).
For citations with binary prestige (Column~2),
$\hat{\beta}_3 = +0.010$ (SE\,=\,0.056, $p\,=\,.863$;
95\,\%\,CI: [$-$0.100, $+$0.120]).
Using the continuous prestige score for downloads (Column~3),
$\hat{\beta}_3 = -0.013$ (SE\,=\,0.032, $p\,=\,.676$;
95\,\%\,CI: [$-$0.076, $+$0.049]).
For citations with continuous prestige (Column~4),
$\hat{\beta}_3 = +0.002$ (SE\,=\,0.028, $p\,=\,.955$;
95\,\%\,CI: [$-$0.053, $+$0.056]).
The confidence intervals span zero in both directions,
indicating the data cannot distinguish the interaction effect from zero.
We find no evidence that AI-style writing modulates the (already
non-significant) institutional prestige association in this context.
In the baseline download model, overall fit is modest ($R^2 = 0.368$),
with publication age the dominant predictor; for citations the fit is
substantially stronger ($R^2 = 0.663$), consistent with the very high
age--citation correlation ($r\,=\,.79$).

\subsection{Robustness Checks}

Table~\ref{tab:robustness} reports six robustness specifications.
Using a binary median-split AI variable (Specification~2) leaves
the interaction non-significant for downloads
($\hat{\beta}_3 = -0.149$, $p=.283$); notably, in this specification
the prestige main effect reaches borderline significance
($\hat{\beta} = 0.159$, $p=.074$), a slight divergence from the
non-significant prestige results in Table~\ref{tab:main} that may
reflect the different covariate structure introduced by the binary AI
operationalization, noting also that the narrow score distribution
(SD\,=\,0.050) means the median split introduces greater information loss
than the continuous specification.
Excluding 2026 articles (Specification~3, $n = 646$) and trimming the
top 5\% of downloads (Specification~4, $n = 800$) yield interaction
estimates within sampling error of the baseline.
Specifications~5 and~6 replicate columns already reported in
Table~\ref{tab:main} (Columns~2 and~3, respectively) and are included
for completeness rather than as independent new checks.
The AI-style main effect on downloads remains positive and significant
across all specifications; all six Prestige\,$\times$\,AI interaction
terms are non-significant.
Full coefficients, $p$-values, and confidence intervals for all
specifications are available in Web Appendix Table~B2.

% ---- 5. Discussion ----------------------------------------------------------
\section{Discussion}
\label{sec:discussion}

\subsection{Summary}

We find that the AI-style writing main effect of \citet{JoLi2026}
replicates for download-based diffusion: formulaic AI-associated
writing patterns independently predict greater download volume on CNKI.
However, the institutional prestige main effect does not replicate:
articles from high-prestige institutions show no significant download
or citation advantage relative to lower-prestige articles in this sample.
As a consequence, the moderation question posed in RQ1 is structurally
moot---one cannot moderate a premium that is not reliably present---and
the interaction is, as expected, null across all four specifications.
Together, these patterns suggest that peer-reviewed journal publication
may constitute a stronger boundary condition than hypothesized, not only
for the AI-moderation mechanism (as anticipated) but also for the
prestige premium itself.

\subsection{Why Is the Prestige Effect Absent?}

The null prestige main effect is theoretically informative.
In the SSRN preprint ecosystem, institutional affiliation functions as
a primary quality heuristic because few alternative signals are
available \citep{JoLi2026}.
In the CNKI journal context, peer review certification provides an
independent quality signal that is correlated with institutional
prestige but not reducible to it \citep{Bornmann2008}.
This pattern is consistent with editorial certification reducing the
marginal role of institutional prestige: once an article has cleared
peer review, readers may weight institutional affiliation less heavily
as a quality signal.
If this mechanism is operative, the AI-style moderation mechanism
becomes structurally moot---an interpretation consistent with, though
not proven by, the present data.
The null result is also consistent with alternative explanations,
including measurement error in the AI-style proxy and insufficient
power at $n=843$; the data do not distinguish between these accounts.

\subsection{Implications for the AI Democratization Thesis}

These findings qualify the claim that AI democratizes academic
knowledge diffusion by homogenizing presentation quality
\citep{Liang2024}.
The equalizing potential documented by \citet{JoLi2026} may be
real in low-gatekeeping, open-access contexts but does not extend---at
least not yet---to formal publication venues with established
quality-signaling mechanisms.
Our results are consistent with editorial certification---rather than
AI-style writing---being the primary factor reducing the marginal
informativeness of institutional affiliation in this context.
Platform designers, editors, and science policymakers should therefore
be cautious about generalizing AI's equalizing potential across
different gatekeeping contexts; the present study's single-context,
observational design means these implications remain questions for
future research rather than confirmed policy recommendations.

\subsection{Limitations}

Several limitations qualify these conclusions.
First, and most critically, the \textit{ai\_style\_score} poses a
\textit{construct validity} threat: as a dictionary-based exact-string
measure of structural writing patterns, it captures AI-\textit{like}
writing style rather than confirmed AI tool use, and may miss
contextual or syntactic AI-associated features not captured by the
23-marker vocabulary.
Authors may adopt formulaic section structure independently of any AI
tool, and conversely AI-assisted abstracts that do not employ
our specific marker phrases will be under-scored.
Furthermore, the marker dictionary was calibrated on 2023--2024
conventions; stylistic drift in 2025--2026 may mean the score
systematically under-represents AI involvement in the most recent
articles, potentially attenuating estimated effects toward zero.
Future work should validate the dictionary-based score against
human annotation or LLM-based evaluation.
Second, the observational OLS design precludes causal inference;
omitted variables (e.g., individual author reputation, reviewer networks)
correlated with both institutional prestige and diffusion outcomes
cannot be ruled out.
Third, the 95\,\% confidence intervals (spanning approximately
$\pm$0.12 log-units for the binary-prestige interaction on downloads,
corresponding to roughly $\pm$0.10 SD of the outcome)
indicate that the present design would be unable to reliably detect
interaction effects smaller than approximately 0.10 standard deviations;
larger samples are needed to rule out small interaction effects with
adequate power.
Fourth, CNKI download counts reflect platform-level behavior including
automated crawlers and bulk-access patterns, not exclusively deliberate
scholarly reading.
Fifth, many articles published in 2025--2026 have yet to accumulate
citations, creating left-censoring that limits the informativeness of
the citation outcome.
Sixth, the sample covers only Chinese-language business and management
articles from a single thematic domain; findings may not generalize
to other disciplines, languages, or national contexts.
These construct-validity limitations are partially offsetting: false
positives (non-AI authors who happen to use the marker phrases) and
false negatives (AI-assisted authors who avoid them) both attenuate
the estimated AI-style effect toward zero---which means the positive
download main effect should be read as a lower bound on the true
association between AI-style writing and diffusion.

% ---- 6. Conclusion ----------------------------------------------------------
\section{Conclusion}
\label{sec:conclusion}

We do not find evidence that AI-style writing attenuates the
institutional reputation premium in CNKI journal diffusion.
More fundamentally, institutional prestige itself is not a robust
predictor of downloads or citations in this peer-reviewed context,
consistent with editorial certification reducing the marginal role
of institutional affiliation in formal publication venues.
This suggests a boundary condition for the mechanism documented by
\citet{JoLi2026}: AI can attenuate reputation-based diffusion only
where a reputation premium is already operative.
Future research should test whether this pattern holds across other
formal publication platforms and disciplines using larger samples, and
should investigate whether direct measures of AI tool use reveal
stronger or qualitatively different patterns.

% ---- References -------------------------------------------------------------
\bibliographystyle{apalike}
\begin{thebibliography}{99}

\bibitem[Biber and Gray(2016)]{Biber2016}
Biber, D., \& Gray, B. (2016).
\newblock \textit{Grammatical complexity in academic English: Linguistic change
  in writing}.
\newblock Cambridge University Press.

\bibitem[Bornmann and Daniel(2008)]{Bornmann2008}
Bornmann, L., \& Daniel, H.-D. (2008).
\newblock What do citation counts measure? A review of studies on citing
  behavior.
\newblock \textit{Journal of Documentation}, 64(1), 45--80.

\bibitem[DiMaggio and Powell(1983)]{DiMaggio1983}
DiMaggio, P., \& Powell, W.~W. (1983).
\newblock The iron cage revisited: Institutional isomorphism and collective
  rationality in organizational fields.
\newblock \textit{American Sociological Review}, 48(2), 147--160.

\bibitem[Jo and Li(2026)]{JoLi2026}
Jo, Y., \& Li, X. (2026).
\newblock Can AI reduce the reputation premium in knowledge diffusion?
  Evidence from SSRN marketing working papers.
\newblock \textit{Marketing Letters} (in press).

\bibitem[Larivi\`{e}re et al.(2010)]{Lariviere2010}
Larivi\`{e}re, V., Macaluso, B., Archambault, \'{E}., \& Gingras, Y. (2010).
\newblock Which scientific elites? On the concentration of research funds,
  publications, and citations.
\newblock \textit{Research Evaluation}, 19(1), 45--53.

\bibitem[Liang et al.(2024)]{Liang2024}
Liang, W., Zhang, Y., Cao, H., Wang, B., Ding, D., Yang, X.,
  \& Zou, J. (2024).
\newblock Mapping the increasing use of LLMs in scientific papers.
\newblock \textit{NEJM AI}, 1(8), AIoa2400196.
\newblock \doi{10.1056/AIoa2400196}

\bibitem[Long et al.(1979)]{Long1979}
Long, J.~S., Allison, P.~D., \& McGinnis, R. (1979).
\newblock Entrance into the academic career.
\newblock \textit{American Sociological Review}, 44(5), 816--830.

\bibitem[MacKinnon and White(1985)]{MacKinnon1985}
MacKinnon, J.~G., \& White, H. (1985).
\newblock Some heteroskedasticity-consistent covariance matrix estimators with
  improved finite sample properties.
\newblock \textit{Journal of Econometrics}, 29(3), 305--325.

\bibitem[Merton(1968)]{Merton1968}
Merton, R.~K. (1968).
\newblock The Matthew effect in science.
\newblock \textit{Science}, 159(3810), 56--63.

\bibitem[Spence(1973)]{Spence1973}
Spence, M. (1973).
\newblock Job market signaling.
\newblock \textit{Quarterly Journal of Economics}, 87(3), 355--374.

\bibitem[Van Noorden(2023)]{VanNoorden2023}
Van Noorden, R. (2023).
\newblock ChatGPT one year on: who is using it, how and why?
\newblock \textit{Nature}, 624(7991), 238--241.

\bibitem[Wang et al.(2013)]{Wang2013}
Wang, D., Song, C., \& Barab\'{a}si, A.-L. (2013).
\newblock Quantifying long-term scientific impact.
\newblock \textit{Science}, 342(6154), 127--132.

\end{thebibliography}

\newpage
% ---- Tables -----------------------------------------------------------------

\begin{table}[htbp]
\centering
\caption{Descriptive Statistics and Bivariate Correlations ($N = 843$)}
\label{tab:descriptives}
\begin{threeparttable}
\footnotesize
\setlength{\tabcolsep}{4pt}
\begin{tabular}{lrrrrr|rrrrrr}
\toprule
& \multicolumn{5}{c|}{\textit{Descriptive statistics}} &
  \multicolumn{6}{c}{\textit{Correlations}} \\
\cmidrule(lr){2-6}\cmidrule(lr){7-12}
Variable & Mean & SD & Min & Mdn & Max &
  (1) & (2) & (3) & (4) & (5) & (6) \\
\midrule
(1) $\log(\text{Downloads}+1)$ & 7.52 & 1.19 & 0.00 & 7.58 & 11.58 &
  ---  &      &     &     &     &    \\
(2) $\log(\text{Citations}+1)$ & 1.81 & 1.52 & 0.00 & 1.61 & 6.96 &
  .47** & --- &     &     &     &   \\
(3) Prestige\_High & 0.668 & 0.471 & 0 & 1 & 1 &
  .05 & .03 & --- &     &     &   \\
(4) ai\_style\_score & 0.049 & 0.050 & 0.000 & 0.044 & 0.261 &
  .05 & $-$.12** & $-$.03 & --- &     &   \\
(5) $\log(\text{Age})$ & 5.69 & 0.94 & 2.08 & 5.85 & 7.12 &
  .56** & .79** & $-$.02 & $-$.10** & --- &  \\
(6) $\log(\text{Abs.\ length})$ & 5.32 & 0.033 & 4.68 & 5.32 & 5.39 &
  .08* & .06  & .04  & .05   & $-$.04 & --- \\
\bottomrule
\end{tabular}
\begin{tablenotes}
\footnotesize
\item \textit{Note.} Prestige\_High\,=\,1 if any author is affiliated with a
  985/C9/211/DFC-ranked institution ($n=563$, 66.8\%).
  ai\_style\_score measures the proportion of 23 pre-specified
  AI-associated structural writing markers present in the Chinese
  abstract (0--1 scale); ai\_z (standardized) is used in regressions.
  $\log(\text{Age})$ = log of days from publication to data collection (June 2026).
  Correlations are Pearson coefficients.
  $^{*}p<.05$;\ $^{**}p<.01$ (two-tailed).
  The small SD for $\log(\text{Abs.\ length})$ (SD\,=\,0.033) is not a
  data entry error: CNKI truncates displayed abstracts at a fixed character
  ceiling ($\approx$\,203 characters), so virtually all abstracts cluster
  at the truncation maximum, producing a right-truncated distribution
  (with most values near the ceiling) and a small standard deviation
  despite a non-trivial minimum.
  The correlation between Prestige\_High and ai\_z ($r=-.03$, $p=.407$)
  confirms the two focal predictors are not collinear.
\end{tablenotes}
\end{threeparttable}
\end{table}

\newpage

\begin{table}[htbp]
\centering
\caption{Main OLS Regression Results: Downloads and Citations ($N = 843$)}
\label{tab:main}
\begin{threeparttable}
\small
\begin{tabular}{lcccc}
\toprule
 & \multicolumn{2}{c}{\textit{Binary prestige (Eq.~1)}}
 & \multicolumn{2}{c}{\textit{Prestige score (Eq.~2)}} \\
\cmidrule(lr){2-3}\cmidrule(lr){4-5}
 & (1) Downloads & (2) Citations & (3) Downloads & (4) Citations \\
\midrule
Prestige\_High
  & $0.113$       & $0.046$       &               &              \\
  & (0.069)       & (0.065)       &               &              \\[4pt]
prestige\_score\_z
  &               &               & $0.031$       & $0.016$      \\
  &               &               & (0.035)       & (0.031)      \\[4pt]
ai\_z
  & $0.216^{***}$ & $0.059$       & $0.187^{***}$ & $0.065^{*}$  \\
  & (0.048)       & (0.043)       & (0.032)       & (0.028)      \\[4pt]
\textbf{Prestige\,$\times$\,ai\_z}
  & $\mathbf{-0.044}$ & $\mathbf{+0.010}$ & $\mathbf{-0.013}$ & $\mathbf{+0.002}$ \\
  & \textbf{(0.062)}  & \textbf{(0.056)}  & \textbf{(0.032)}  & \textbf{(0.028)}  \\
  & \textbf{[$-$.164, $+$.077]} & \textbf{[$-$.100, $+$.120]}
  & \textbf{[$-$.076, $+$.049]} & \textbf{[$-$.053, $+$.056]} \\[4pt]
$\log(\text{Age})$
  & $0.697^{***}$ & $0.900^{***}$ & $0.702^{***}$ & $0.901^{***}$ \\
  & (0.090)       & (0.093)       & (0.090)       & (0.093)       \\[4pt]
$\log(\text{Abs.\ length})$
  & $1.825$       & $1.395$       & $1.876$       & $1.404$       \\
  & (1.601)       & (2.224)       & (1.629)       & (2.228)       \\[4pt]
\midrule
Year FEs           & Yes  & Yes  & Yes  & Yes  \\
$R^2$              & .368 & .663 & .366 & .663 \\
Adj.\ $R^2$        & .362 & .660 & .360 & .660 \\
$N$                & 843  & 843  & 843  & 843  \\
\bottomrule
\end{tabular}
\begin{tablenotes}
\footnotesize
\item \textit{Note.} OLS with HC3 heteroskedasticity-robust standard errors
  in parentheses. 95\,\% confidence intervals for the interaction term in
  brackets. The interaction coefficient ($\hat{\beta}_3$) is the parameter of
  central inferential interest and is shown in \textbf{bold}.
  Columns 1--2 use binary Prestige\_High (Equation~1);
  Columns 3--4 use standardized four-tier prestige score (Equation~2).
  Intercept and year dummy coefficients omitted for brevity; available in
  Web Appendix Table~B1.
  $^{*}p<.05$;\ $^{**}p<.01$;\ $^{***}p<.001$.
\end{tablenotes}
\end{threeparttable}
\end{table}

\newpage

\begin{table}[htbp]
\centering
\caption{Robustness Checks: Key Coefficients Across Alternative Specifications}
\label{tab:robustness}
\begin{threeparttable}
\small
\begin{tabular}{llp{4.2cm}ccc}
\toprule
 & & & \multicolumn{3}{c}{\textit{Focal coefficients}} \\
\cmidrule(lr){4-6}
Spec. & Outcome & Modification
  & Prestige & AI-Style & \textbf{Prestige\,$\times$\,AI} \\
\midrule
1 & Downloads & Baseline (binary prestige, Eq.~1)
  & $0.113$ & $0.216^{***}$  & $-0.044$ \\
  &           &
  & (0.069) & (0.048)        & (0.062)  \\[4pt]
2 & Downloads & Binary AI (median split) instead of continuous
  & $0.159^{\dagger}$ & $0.383^{***}$ & $-0.149$  \\
  &           &
  & (0.089)           & (0.108)       & (0.139)  \\[4pt]
3 & Downloads & Exclude 2026 articles ($n = 646$)
  & $0.097$ & $0.174^{***}$  & $-0.001$ \\
  &           &
  & (0.073) & (0.053)        & (0.068)  \\[4pt]
4 & Downloads & Trim top 5\% download outliers ($n=800$)
  & $0.074$ & $0.213^{***}$  & $-0.042$ \\
  &           &
  & (0.065) & (0.046)        & (0.059)  \\[4pt]
5 & Citations & Baseline (binary prestige, Eq.~1)
  & $0.046$  & $0.059$  & $+0.010$ \\
  &           &
  & (0.065)  & (0.043)  & (0.056)  \\[4pt]
6 & Downloads & Continuous prestige score (Eq.~2)
  & $0.031$ & $0.187^{***}$ & $-0.013$ \\
  &           &
  & (0.035) & (0.032)       & (0.032)  \\[4pt]
\bottomrule
\end{tabular}
\begin{tablenotes}
\footnotesize
\item \textit{Note.} Each row is an independent OLS regression.
  Cells report $\hat{\beta}$ with HC3 robust SE in parentheses.
  All controls (log age, log abstract length, year FEs)
  are included in every specification.
  The Prestige\,$\times$\,AI interaction (bold column) is not statistically
  significant in any of the six specifications.
  Full regression tables are available in Web Appendix Table~B2.
  $^{\dagger}p<.10$;\ $^{*}p<.05$;\ $^{**}p<.01$;\ $^{***}p<.001$.
\end{tablenotes}
\end{threeparttable}
\end{table}

\newpage
% ---- Web Appendix Outline ---------------------------------------------------
\section*{Web Appendix (not for print)}

\subsection*{Appendix A: Data Construction and Variable Coding}

\textbf{A1. Sample construction.}
Full CNKI search string, inclusion and exclusion criteria, and
article-count flowchart by year and topic domain.

\textbf{A2. AI-style score: construction and validation.}
Complete list of all 23 Chinese structural and formulaic markers
(section-header phrases and boilerplate templates) with Chinese
original text, English translation, and example abstracts;
description of exact string-matching scoring procedure;
score distribution histogram and percentile table;
inter-coder agreement for dictionary construction ($\kappa=0.79$,
development sample).

\textbf{A3. Institution prestige coding.}
Full rubric and list of C9, 985, and 211/DFC institutions;
prestige-score distribution by year; inter-rater reliability ($\kappa=0.93$).
Sensitivity analysis using first-author institution prestige coding;
focal regression coefficients compared with maximum-tier baseline.

\textbf{A4. Variable dictionary.}
Definitions, measurement units, and additional descriptive statistics
for all variables, disaggregated by year and prestige tier.

\subsection*{Appendix B: Additional Analyses}

\textbf{B1. Full regression tables.}
Tables~\ref{tab:main} and \ref{tab:robustness} with intercept,
all year dummy coefficients, and all control coefficients.

\textbf{B2. Full robustness regression tables.}
All coefficients for all six robustness specifications in
Table~\ref{tab:robustness}.

\textbf{B3. Marginal effects plots.}
Predicted $\log(\text{Downloads}+1)$ and $\log(\text{Citations}+1)$
as a function of \textit{ai\_z}, plotted separately for high- and
low-prestige groups, illustrating the near-parallel slopes corresponding
to the non-significant interaction.

\textbf{B4. AI-style score distribution by prestige tier.}
Bar chart of mean AI-style scores across the four prestige tiers (1--4),
with group-level descriptive statistics.

\textbf{B5. Negative-binomial models on raw counts.}
Parallel models estimated on raw (untransformed) download and citation
counts using negative-binomial regression; coefficients and incidence
rate ratios.

\textbf{B6. Topic fixed-effect sensitivity.}
Results adding topic/theme binary dummies (AI, supply chain, ESG,
governance, innovation) to the main specifications.

\textbf{B7. Marketing and management sub-sample sensitivity.}
Main regression results restricted to articles from CNKI journals
classified in the marketing or business administration sub-categories,
confirming that findings are not driven by non-business topic domains.

\subsection*{Appendix C: Comparison with Jo and Li (2026)}

Side-by-side coefficient table comparing focal estimates from
\citet{JoLi2026} (SSRN, marketing working papers) with estimates from
the present study (CNKI, peer-reviewed Chinese journal articles),
with narrative discussion of cross-context differences in sample
composition, platform mechanics, language, and prestige-measurement
approach.

\end{document}
